\documentclass[11pt]{article}
\usepackage[margin=25mm]{geometry}
\usepackage{amsmath,amssymb,amsthm,booktabs}
\usepackage[colorlinks=true,urlcolor=blue,linkcolor=blue]{hyperref}
\usepackage[T1]{fontenc}
\newtheorem{theorem}{Theorem}
\newtheorem{proposition}{Proposition}
\setlength{\parindent}{0pt}
\setlength{\parskip}{6pt}
\title{Support projection can destroy backward-error feasibility\\\large Conjecture 00000005137}
\author{zbw123456}\date{14 October 2026}
\begin{document}\maketitle
\textbf{Scope.} We test the stated support-projection rule for matrix-only
backward error in $Ax=b$, using the Frobenius norm and coordinate projection
onto the nonzero positions of $A$. These choices are stated explicitly because
the conjecture does not specify a norm or which input data can be perturbed.

Take
\[
 A=\begin{pmatrix}1&1\\0&1\end{pmatrix},\quad
 x=\binom11,\quad b=\binom32,\quad r=b-Ax=\binom11.
\]
The support $S=\{(1,1),(1,2),(2,2)\}$ is a tree in the bipartite row-column
graph: its path is column 1, row 1, column 2, row 2.
For $E=\left(\begin{smallmatrix}a&b'\\c&d\end{smallmatrix}\right)$, the backward
error constraint is
\[
 (A+E)x=b\quad\Longleftrightarrow\quad a+b'=1,\quad c+d=1.
\]
Define
\[
 D=\begin{pmatrix}1/2&1/2\\1/2&1/2\end{pmatrix},\qquad
 T=\begin{pmatrix}1/2&1/2\\0&1\end{pmatrix}.
\]
\begin{theorem}
$D$ is the unique minimum-Frobenius-norm dense feasible perturbation, with
$\|D\|_F^2=1$. The minimum supported perturbation is $T$, with
$\|T\|_F^2=3/2$. The support projection of $D$ is not feasible.
\end{theorem}
\begin{proof}
For every feasible $E$,
\[
 \|E\|_F^2
 =1+\tfrac12(a-b')^2+\tfrac12(c-d)^2\ge1.
\]
Equality requires $a=b'$ and $c=d$, so the row sums force all four entries
to equal $1/2$. Under the support constraint $c=0$, feasibility forces $d=1$,
and
\[
 \|E\|_F^2=\tfrac32+\tfrac12(a-b')^2\ge\tfrac32,
\]
with equality at $T$. But
\[
 \Pi_S(D)=\begin{pmatrix}1/2&1/2\\0&1/2\end{pmatrix},\qquad
 \Pi_S(D)x=\binom{1}{1/2}\ne r.
\]
Therefore projection of the dense minimizer cannot give a sparse minimizer.
There is also a strict discrepancy between the two optimal squared norms.
\end{proof}

\textbf{Formal correspondence.} \texttt{SparseBackward.lean} uses four real
coordinates for an actual $2\times2$ matrix, explicit matrix-vector multiplication,
and the actual equation $(A+E)x=b$. Its minimization predicates quantify over
all four real entries of every competitor. The squared objective is the squared
Frobenius norm, so it has exactly the same minimizers. Lean proves dense
optimality, sparse optimality, uniqueness of the dense minimizer, and failure
of feasibility after projection. In particular,
\texttt{no\_projected\_dense\_minimizer} excludes every dense minimizer, not just
a selected feasible perturbation. \texttt{SupportTree.lean} proves Mathlib's connected-and-acyclic tree predicate
for the four-vertex path. \texttt{support\_correspondence} identifies its
row-column adjacencies with the actual nonzero entries of $A$.
\texttt{counterexample\_with\_tree} combines this verified tree support
with the minimization counterexample.

\textbf{Interpretation limit.} A different norm, perturbations of $b$ as well,
or a support-aware correction performed \emph{after} projection would define a
different problem. This result refutes direct coordinate projection under the
stated model. The competition must decide whether that is the intended meaning
of its unspecified conversion formula; no acceptance is claimed.

\textbf{Reproduction and sources.} Lean 4.19.0 and pinned Mathlib are specified
in the project. The final axiom audit contains only \texttt{propext},
\texttt{Classical.choice}, and \texttt{Quot.sound}. The official bilingual text
is included. For the distinction between structured and unstructured backward
error, see D.~J. Higham and N.~J. Higham,
\href{https://eprints.maths.manchester.ac.uk/347/}{Backward error and condition of structured linear systems} (1992).
\end{document}
